\section{Conclusioni}
\label{sec:conclusions}

Rispondo esplicitamente alle tre domande richieste dalla consegna.

\subsection{Cosa non ha funzionato}
\begin{itemize}
  \item \textbf{Acrobot non decolla}: plateau a $\approx -345$ contro un ottimo di
        $\approx -100$, e nessun miglioramento per 15\,000 episodi. Raddoppiare il
        budget non serve, perch\'e il limite \`e la risoluzione della Q-table
        (Sez.~\ref{sec:res-acrobot-control}).
  \item \textbf{CartPole non raggiunge 500}: con $10^4$ stati il tetto in 10\,000
        episodi \`e circa met\`a del massimo. Anche qui il vincolo \`e rappresentativo,
        non algoritmico.
  \item \textbf{I tre fallimenti indotti} (Sez.~\ref{sec:failure}): $\alpha=1.0$ non
        converge, $\varepsilon=0$ resta subottimale, $\gamma=0$ non pu\`o imparare
        affatto --- e in tutti e tre il codice gira correttamente.
  \item \textbf{Double DQN non supera il DQN base} a 200k step, anzi lo peggiora: il
        maximization bias che dovrebbe correggere non ha ancora avuto il tempo di
        accumularsi (Sez.~\ref{sec:res-pong}).
  \item \textbf{Il DQN non era riproducibile}: lo spazio delle azioni ha un RNG separato
        che n\'e \code{set\_seed} n\'e \code{reset(seed=)} toccano, e da l\`i viene ogni
        azione esplorativa; due run con lo stesso seed divergevano. Scoperto e corretto
        (Sez.~\ref{sec:discussion}).
  \item \textbf{Il mio primo confronto MC vs TD(0) non misurava nulla}: graficava il
        return della policy condivisa, identico per costruzione tra i due stimatori ---
        ``giusto'' e inutile (Sez.~\ref{sec:res-cartpole-pred}).
\end{itemize}

\subsection{Cosa mi ha sorpreso}
\begin{itemize}
  \item \textbf{Le due ipotesi smentite su Acrobot}: pi\`u bin sono \emph{meglio} (non
        peggio --- l'aliasing di una griglia grossolana domina la curse of dimensionality),
        e \emph{meno} esplorazione \`e meglio a $\varepsilon$ costante, salvo poi essere
        battuta dall'annealing: il problema non era ``quanta'' esplorazione ma ``quando''
        (Sez.~\ref{sec:res-acrobot-bins}, \ref{sec:res-acrobot-eps}).
  \item \textbf{SARSA batte Q\hbox{-}Learning su CartPole}, con varianza minore. Mi
        aspettavo il contrario; ma la curva misura la performance \emph{con
        l'esplorazione attiva}, e SARSA \`e l'unico dei due che ne tiene conto.
  \item \textbf{Il bias di TD(0) non \`e una costante}: sulla policy casuale (episodi
        corti) TD raggiunge MC, su quella appresa ($\approx 330$ passi) resta indietro del
        $13\%$ --- dipende dalla distanza temporale dal reward, che il \emph{random walk}
        del corso nasconde.
  \item \textbf{Senza target network l'apprendimento \`e una lotteria}: mi aspettavo un
        agente ``pi\`u lento'', ho trovato invece che 2 seed su 3 non decollano affatto e
        solo uno impara. La media ($-20$) nasconde il fenomeno, il per-seed lo rivela: la
        target network non accelera soltanto, rende l'apprendimento \emph{affidabile}
        (Sez.~\ref{sec:res-pong}).
\end{itemize}

\subsection{Limiti degli algoritmi}
\begin{itemize}
  \item \textbf{Discretizzazione} --- il collo di bottiglia dominante, e ha due lame:
        pochi bin non distinguono stati che richiedono azioni diverse, troppi bin fanno
        esplodere lo spazio. Nessuna scelta di bin risolve entrambi i problemi: serve
        l'approssimazione di funzione. \`E la motivazione empirica del DQN.
  \item \textbf{Efficienza campionaria e convergenza} --- migliaia di episodi per
        risultati modesti (nessuna generalizzazione tra stati simili), e con $\alpha$ ed
        $\varepsilon$ costanti la stima fluttua senza mai convergere davvero.
  \item \textbf{Sparsit\`a del reward} --- su Acrobot il segnale ($-1$ per passo) non
        dice \emph{quale} azione fosse sbagliata: il credit assignment su 500 passi \`e il
        problema vero, e nessuno dei quattro algoritmi tabellari lo affronta.
  \item \textbf{Sensibilit\`a agli iperparametri} --- marcata, e in modi che cambiano da
        ambiente ad ambiente e persino da algoritmo ad algoritmo ($\alpha$ alto danneggia
        SARSA ma non Q\hbox{-}Learning).
  \item \textbf{Costo del Deep RL} --- il DQN risolve la rappresentazione, ma 200k step
        sono lo $0.4\%$ del budget Nature: \`e per questo che Pong si ferma a
        $\approx -13$, e che senza target network met\`a dei seed non decolla.
\end{itemize}

